% Part: applied-modal-logics
% Chapter: temporal-logic
% Section: introduction

\documentclass[../../../include/open-logic-section]{subfiles}

\begin{document}

\olfileid{aml}{tl}{int}

\olsection{प्रस्तावना}

कालिक तर्कशास्त्रे यापुढे काय घडेल किंवा यापूर्वी
काय घडले आहे यांविषयीच्या विधानांचा विचार करतात.
आर्थर प्रायर यांना कालिक तर्कशास्त्राचे जनक मानले
जाते; त्यांनी त्याला \emph{काळवाचक तर्कशास्त्र}
असे नाव दिले. विधानीय तर्कशास्त्राची मूलभूत चौकट
विधाने साधारणतः काळवाचक भेद नसलेली मानते. त्या
चौकटीत कालिक संकारक घालण्याच्या प्रायर यांच्या
मूळ मोडल पद्धतीचे येथे आपण मुख्यतः अनुसरण करू.

उदाहरणार्थ, विधानीय तर्कशास्त्रात मी बीझी नावाच्या
कुत्र्याबद्दल बोलू शकतो. कुत्र्यांच्या स्वभावाप्रमाणे
ते कधी बसते आणि कधी बसत नाही. बीझी बसले आहे
आणि बीझी बसलेले नाही असे एकत्र म्हणणे अभिजात
तर्कशास्त्रात विरोधाभासी ठरेल. पण ही दोन्ही विधाने
खरी असू शकतात, फक्त एकाच वेळी नाहीत. भाषेत
कालिक संकारक घातल्यास ही गोष्ट तुलनेने सहज
व्यक्त करता येते. कालिक संकारकांमुळे ``बीझीला
खाऊ किंवा चेंडू मिळेल'' यावरून ``बीझीला खाऊ
मिळेल किंवा बीझीला चेंडू मिळेल'' यांसारख्या
अनुमानांची वैधताही स्पष्ट करता येते.

मात्र कालिक तर्कशास्त्रामुळे अनेक तत्त्वज्ञानविषयक
प्रश्न निर्माण होतात आणि त्यामुळे त्याची एक चौकट
निवडण्याऐवजी दुसरी निवडावी लागू शकते. उदाहरणार्थ,
भविष्यविषयक परिस्थितीवश विधान म्हणजे भविष्याबद्दलचे
असे विधान की जे अनिवार्यही नाही आणि अशक्यही नाही.
``रिचर्ड उद्या किराणा दुकानात जाईल'' असे म्हटल्यास,
आपण अजून न घडलेल्या घटनेविषयी विधान करतो; त्या
विधानाचे सत्यतामूल्य वादग्रस्त असते. किंबहुना,
त्याला मुळात सत्यतामूल्य \emph{देता} तरी येते का,
हाही वादाचा विषय आहे. आपण कठोर निर्धारणवादी
असू, तर संबंधित घटना घडण्यापूर्वीच हे विधान
खरे किंवा खोटे असते, एवढेच की त्याचे सत्यतामूल्य
आपल्याला अजून माहीत नसते, असे मानणे आपल्याला
मान्य होऊ शकते. याउलट, भविष्य खऱ्या अर्थाने
खुले आहे आणि भविष्यविषयक परिस्थितीवश विधानांची
सत्यतामूल्ये अजून अनिर्धारित आहेत, असेही आपण
मानू शकतो.

काळाची रचना आणि स्वरूप यांविषयीच्या अशा अनेक
भूमिका कालिक तर्कशास्त्रासाठी आपण निवडलेल्या
प्रतिमानांत आणि सैद्धांतिक चौकटींत अंतर्भूत असतात.
उदाहरणार्थ, काळ रेषीय, शाखायुक्त किंवा अगदी
वर्तुळाकार आहे अशी प्रतिमाने तयार करावीत का,
हा प्रश्न आपण विचारू शकतो. कालिक प्रतिमानांना
आरंभबिंदू आणि अंतबिंदू असावेत का, तसेच काळाचे
निरूपण विविक्त क्षणांनी करावे की सांतत्य म्हणून
करावे, हेही ठरवावे लागेल.

\end{document}
